\section{Ruta de lectura con la biblioteca local y compartida}
\label{sec:ruta-biblioteca-local-drive}

Esta sección convierte las bibliotecas señaladas por la autora en un programa
de estudio verificable. No se infiere el contenido de un libro a partir de su
nombre: se revisaron portadas, páginas legales, índices y texto extraíble de
los tres ejemplares de cálculo fraccionario disponibles localmente, y se
inspeccionaron los índices de los tres volúmenes directamente pertinentes de
la carpeta compartida de Google Drive. La finalidad no es acumular citas, sino
asignar a cada fuente una pregunta matemática concreta.

\subsection{Un solo orden de lectura}
\label{sec:ruta-unica-estudio}

Ésta es la ruta oficial del volumen. Las actividades situadas en otros
capítulos sirven para practicar una etapa, pero no abren itinerarios
paralelos. Se avanza de manera acumulativa: una etapa se da por terminada
cuando puede realizarse su comprobación de salida sin copiar fórmulas.

\begin{longtable}{@{}L{1.1cm}L{2.4cm}L{5.5cm}L{5.05cm}@{}}
\caption{Orden progresivo antes, durante y después del contenido de tesis.}
\label{tab:ruta-unica-biblioteca}\\
\toprule
Etapa & Momento & Lectura y archivos & Comprobación para avanzar \\
\midrule
\endfirsthead
\toprule
Etapa & Momento & Lectura y archivos & Comprobación para avanzar \\
\midrule
\endhead
0 & Preliminar opcional &
Diagnóstico con \path{[Anton] Algebra lineal.pdf},
\path{[Marsden] Calculo Vectorial 3rd Ed.PDF} y
\texttt{Cálculo Infinitesimal - Michael Spivak - Universidad de Brandeis.pdf},
en la subcarpeta \texttt{Cálculo vectorial y álgebra lineal}. &
Calcular determinantes, valores propios, derivadas parciales, un jacobiano y
una integral impropia. Si estos productos ya son fluidos, no se releen los
libros completos. \\
1 & Antes: base analítica &
Kantorovich, archivo \path{978-3-030-91222-2.pdf}: derivadas,
integrales, varias variables, series y EDO, pp.\ 243--760
\cite{Kantorovich2022Fundamentals}. Oliveira: Gamma, beta,
Mittag--Leffler y transformadas \cite{Oliveira2019SolvedFractional}. &
Derivar una identidad beta--Gamma, resolver un sistema lineal por resolvente y
justificar las condiciones de una transformada de Laplace. \\
2 & Antes: EDO y topología &
Los tres libros UNAM; Hirsch--Smale--Devaney, archivo cuyo nombre comienza
\path{[Pure and Applied Mathematics (Academic Press), 60] Hirsch...pdf},
caps.\ 1--10 \cite{HirschSmaleDevaney2013}. &
Distinguir espacio, métrica y topología; demostrar existencia local bajo las
hipótesis declaradas; definir flujo, órbita, invariancia, conjunto
\(\omega\)-límite y cuenca. \\
3 & Antes: dinámica y caos entero &
 Hirsch--Smale--Devaney, caps.\ 12 y 14--16; Hoppensteadt,
 \path{(Applied Mathematical Sciences, 94) Frank C. Hoppensteadt - Analysis and Simulation of Chaotic Systems-Springer (2000).pdf};
 Wimberger, \path{978-3-031-01249-5.pdf}, pp.\ 9--152
 \cite{Wimberger2022Nonlinear}; Vallejo--Sanjuán,
 \path{978-3-319-51893-0.pdf}; y los capítulos 5--6 de
 \path{978-3-662-53685-8.pdf}. &
Construir un mapa de Poincaré sencillo; separar ciclo, bifurcación,
homoclinicidad, herradura, exponente finito y definición topológica de caos. \\
4 & Antes: cálculo fraccionario &
 Oliveira, Herrmann y Das, en
 la colección local identificada en este reporte como
 \path{Biblioteca local/Calculo Fraccionario}; después, la
 \cref{sec:geometria-topologia-caputo}; como contraste numérico especializado,
 \path{978-981-19-3273-1.pdf}. &
Pasar de Caputo a Volterra; declarar terminal inferior e historia; refutar la
ley de semigrupo puntual con Mittag--Leffler; evaluar
\((\mathrm i\omega)^q\) en una rama declarada. \\
5 & Durante: métodos de localización &
Kantorovich II, \path{978-3-031-46320-4.pdf}: álgebra lineal,
análisis complejo, Fourier y Laplace, pp.\ 1--717
\cite{Kantorovich2024Advanced}; Cicogna,
\path{978-3-030-59472-5.pdf}, pp.\ 57--127
\cite{Cicogna2020Exercises};
\cref{sec:lure_transferencia,sec:machado-fdf,sec:integracion_continuacion,sec:perpetual-points}. &
Reconstruir sin saltos una transferencia, una función descriptiva, una
semilla modal y una homotopía; explicar por qué Machado y puntos perpetuos son
localizadores complementarios y no pruebas de ocultedad. \\
6 & Durante: protocolos y casos &
\cref{sec:modelos,sec:fundamentos}; expedientes enteros y fraccionarios del
volumen; Pham--Volos--Kapitaniak,
\path{978-3-319-53721-4.pdf}; Wang--Kuznetsov--Chen,
\path{978-3-030-75821-9.pdf}; volumen de Pham et al.\ con nombre largo
\path{(Studies in Systems, Decision and Control 133)...Hidden Attractor.pdf}. &
Para un caso, entregar modelo, parámetros, equilibrios, estabilidad, método
de semilla, contrato de memoria, clasificación y sondeo de cuenca, indicando
qué conclusión permanece abierta. \\
7 & Después: profundización &
 Kuznetsov, \path{978-3-031-22007-4.pdf}; Amster; teoría geométrica II;
 Putkaradze, \path{978-3-031-84977-0.pdf}, caps.\ 3--9; el capítulo
 geométrico de Lorenz en \path{978-3-0348-0697-8.pdf}; y material posterior
 de la \cref{sec:geo-study-plan}. &
Formular una continuación de bifurcaciones, una pregunta de grado o punto
fijo y un protocolo geométrico que distinga evidencia numérica de teorema. \\
8 & Futuro: extensión de tesis &
\cref{sec:candidate-selection,sec:cand-munoz,sec:cand-jerk,sec:cand-extended-liu};
capítulos selectivos de circuitos y control en
\path{978-3-031-55460-5.pdf} y \path{978-3-031-36364-1.pdf}. &
Proponer un candidato con ecuaciones congeladas, prueba de compatibilidad
estructural, oráculo numérico, criterio de continuación y ensayo de cuenca
predeclarado. \\
\bottomrule
\end{longtable}

\begin{remark}[cómo usar la ruta]
No se recomienda imprimir ni leer linealmente todos los PDF de la biblioteca.
Los textos núcleo sostienen las etapas; los demás se consultan por capítulos
cuando una pregunta de tesis los activa. Cada etapa produce una demostración,
un cálculo o un contrato verificable, no sólo páginas leídas.
\end{remark}

\subsection{Proveniencia y límite de la inspección}

El primer corpus está en
\path{Biblioteca local/Calculo Fraccionario} y contiene los ejemplares de
Oliveira, Herrmann y Das. El segundo se consultó en la carpeta compartida
\url{https://drive.google.com/drive/folders/12G8I9Of-uC4HMMmruElm3pGb7UzYvN9L}.
Allí se verificaron, entre otros, Hirsch--Smale--Devaney,
Hoppensteadt y el volumen colectivo editado por Pham, Vaidyanathan, Volos y
Kapitaniak \cite{HirschSmaleDevaney2013,Hoppensteadt2000Chaotic,
PhamEtAl2018HiddenAttractors}. La vista de Drive confirmó cincuenta elementos
en la raíz y permitió inspeccionar además la subcarpeta
\texttt{Cálculo vectorial y álgebra lineal} y el volumen
\path{Introducing Fractals - Nigel Lesmoir-Gordon.pdf}. Los archivos cuyo
nombre es sólo un eISBN se resolvieron contra fichas editoriales primarias
cuando podían aportar prerrequisitos, dinámica, transformadas,
bifurcaciones, atractores ocultos o realización circuital. Los restantes
permanecen inventariados sin convertirlos automáticamente en lecturas de
tesis. Esta regla evita atribuir a un archivo una edición o un tema
equivocados.

La ruta original de la unidad \path{G:} no fue legible desde el entorno de
compilación. El enlace de Drive permitió, sin embargo, inspeccionar la carpeta
compartida que la usuaria indicó como su segunda biblioteca. La inspección
directa de un índice o de una ficha editorial se distingue de la simple
confirmación de que el nombre del archivo existe.

\subsection{Tres libros de cálculo fraccionario, tres funciones distintas}

\begin{table}[H]
\centering
\caption{Función pedagógica de los libros fraccionarios inspeccionados.}
\label{tab:ruta-frac-books}
\small
\begin{tabularx}{\textwidth}{@{}L{0.24\textwidth}L{0.29\textwidth}Y@{}}
\toprule
Fuente & Núcleo verificado & Papel dentro de esta metodología \\
\midrule
Oliveira \cite{Oliveira2019SolvedFractional} &
Funciones especiales (pp.\ 17--67), Mittag--Leffler (69--113), transformadas
integrales (115--167), derivadas fraccionarias (169--222) y aplicaciones
(223--307). &
Es la fuente de trabajo algebraico: Gamma, Mittag--Leffler, Laplace,
funciones de Green, ecuaciones integrales de Volterra y ejercicios con
solución. \\
Herrmann \cite{Herrmann2014Fractional} &
Operadores fraccionarios y semigrupo en el orden (caps.\ 3 y 5), oscilador
fraccionario (cap.\ 6), no localidad y memoria (cap.\ 8), y formalismos de
simetría, tensores y campos (caps.\ 13, 16 y 23--25). &
Ofrece intuición física y puentes geométricos. Sus ``grupos fraccionarios'' y
``fractors'' son formalismos específicos; no se presentan como geometría
universal de sistemas Caputo. \\
Das \cite{Das2020Kindergarten} &
Integración (pp.\ 47--99), derivadas (100--154), inicialización y Laplace
generalizada (210--254), ecuaciones diferenciales e integrales fraccionarias
(255--343) y ramas complejas en el apéndice E (466--470). &
Es un puente conceptual amplio, útil para terminal inferior, inicialización,
sistemas vectoriales y ramas de \(s^q\). Su tono deliberadamente informal
requiere respaldar los teoremas con una fuente analítica rigurosa. \\
\bottomrule
\end{tabularx}
\end{table}

Los números anteriores son páginas impresas. Para localizar rápidamente el
material en los PDF locales se usa la correspondencia de la
\cref{tab:ruta-frac-pdf-pages}.

\begin{longtable}{@{}L{0.21\textwidth}L{0.35\textwidth}L{0.35\textwidth}@{}}
\caption{Correspondencia entre tema, página impresa y página PDF.}
\label{tab:ruta-frac-pdf-pages}\\
\toprule
Fuente & Tema y página impresa & Página del PDF \\
\midrule
\endfirsthead
\toprule
Fuente & Tema y página impresa & Página del PDF \\
\midrule
\endhead
Oliveira & Mittag--Leffler, pp.\ 69--113 & PDF 81--125 \\
Oliveira & Transformadas, pp.\ 115--167 & PDF 126--178 \\
Oliveira & Derivadas fraccionarias, pp.\ 169--222 & PDF 179--232 \\
Oliveira & Volterra/Green: ejercicios 33, 35, 36 y 38,
pp.\ 277, 280, 281 y 283 & PDF 287, 290, 291 y 293 \\
Das & Riemann--Liouville y Caputo, \S\S3.7--3.10,
pp.\ 119--125 & PDF 156--162 \\
Das & Inicialización, cap.\ 5, pp.\ 210--254 & PDF 247--291 \\
Das & Sistemas vectoriales, \S7.13, p.\ 322 y \S7.17, p.\ 335 &
PDF 359 y 372 \\
Das & Ramas complejas y superficies de Riemann, apéndice E,
pp.\ 466--470 & PDF 503--507 \\
Herrmann & Definiciones y operadores, \S3.4.2 y \S\S5.2--5.3,
pp.\ 26 y 45--52 & PDF 47 y 66--73 \\
Herrmann & Interpretación geométrica elemental, \S5.7, p.\ 62 & PDF 83 \\
Herrmann & Semigrupo en el orden, \S5.9.1, p.\ 66 & PDF 87 \\
Herrmann & No localidad y memoria, cap.\ 8, pp.\ 87--106 & PDF 108--127 \\
Herrmann & Tensores, campos y gauge, caps.\ 23--25,
pp.\ 353--378 & PDF 374--399 \\
\bottomrule
\end{longtable}

\subsection{La distinción que organiza toda la lectura fraccionaria}

Las fuentes contienen identidades de composición de integrales, por ejemplo
\begin{equation}
 I^\alpha I^\beta=I^{\alpha+\beta}.
 \label{eq:library-order-semigroup}
\end{equation}
La igualdad es un semigrupo respecto al \emph{orden} del operador. No implica
la ley temporal
\begin{equation}
 \Phi(t+s,x_0)=\Phi\bigl(t,\Phi(s,x_0)\bigr)
 \label{eq:library-false-time-semigroup}
\end{equation}
para el estado instantáneo de un PVI Caputo. La
\cref{sec:geometria-topologia-caputo} deriva la razón desde la ecuación de
Volterra: al desplazar el tiempo aparece un perfil de memoria que no queda
determinado sólo por \(x(s)\). Éste es el punto donde el cálculo fraccionario
se convierte en una pregunta de sistemas dinámicos sobre un espacio de
historias.

Análogamente, una fórmula generalizada de Leibniz o de cadena debe conservarse
con la definición exacta bajo la cual fue obtenida. No se puede sustituir por
la regla ordinaria
\[
 D_C^q[f(x(t))]=Df(x(t))D_C^q x(t),
\]
que es falsa en general. Los desarrollos de puntos perpetuos y de cambio de
coordenadas del informe se construyen precisamente sin usar esa identidad.

\subsection{Ramas complejas y función de transferencia fraccionaria}

La transferencia fraccionaria evalúa \(z=s^q\). En el eje imaginario,
\begin{equation}
 (\ii\omega)^q
 =\omega^q\exp\!\left(\ii q\frac{\pi}{2}\right),
 \qquad \omega>0,
 \label{eq:library-principal-branch}
\end{equation}
si se declara la rama principal. El apéndice de Das sobre ramas y superficies
de Riemann ofrece un apoyo visual para entender por qué \(s^q\) es multivaluada
antes de elegir una rama. Esta topología pertenece al análisis complejo de la
transferencia; no es la topología de una cuenca o de un atractor. Mantener
separadas ambas nociones evita usar la palabra ``topológico'' sin especificar
el objeto.

El control mínimo para todo cálculo de frecuencia es, por tanto:
\begin{enumerate}
 \item declarar el corte y la rama de \(s^q\);
 \item comprobar \(q=1\) como límite entero;
 \item verificar por separado partes real e imaginaria;
 \item no cambiar la rama durante una continuación;
 \item registrar la convención de realimentación de \(W_+\) o \(W_-\).
\end{enumerate}

\subsection{La biblioteca geométrica y caótica de Drive}

El índice de Hirsch--Smale--Devaney
\cite{HirschSmaleDevaney2013} permite una ruta especialmente cercana al
proyecto:
\begin{enumerate}
 \item capítulos 1--6: ecuaciones y sistemas lineales, retratos y
 clasificación;
 \item capítulos 7--9: existencia, dependencia continua, ecuación variacional,
 equilibrios y técnicas globales;
 \item capítulo 10: secciones locales, mapa de Poincaré, conjuntos límite y
 Poincaré--Bendixson;
 \item capítulo 12: circuitos, Liénard, Van der Pol y bifurcación de Hopf;
 \item capítulo 14: sistema y atractor de Lorenz;
 \item capítulo 15: dinámica discreta, caos, dinámica simbólica, corrimiento y
 conjunto de Cantor;
 \item capítulo 16: Shilnikov, herradura, doble scroll, homoclinicidad y una
 exploración del circuito de Chua;
 \item capítulo 17: prueba de existencia, unicidad, dependencia continua y
 diferenciabilidad del flujo.
\end{enumerate}

Hoppensteadt \cite{Hoppensteadt2000Chaotic} añade una ruta aplicada que encaja
con dos expedientes concretos: ecuaciones angulares y cilindro del PLL en
\S2.2; sistemas disipativos y PLL en \S2.4; mapas estroboscópicos en \S2.5;
variedades estable/inestable en \S3.2.3; estabilidad orbital y retorno de
Poincaré en \S3.5; y bifurcación, función implícita y métodos topológicos en el
capítulo 4. Así, la geometría cilíndrica del PLL deja de ser una corrección de
dibujo y se convierte en parte del modelo matemático.

El volumen de Pham et al. \cite{PhamEtAl2018HiddenAttractors} cumple otra
función. Su índice separa una parte de sistemas con atractores autoexcitados y
otra de sistemas con atractores ocultos. La segunda incluye órbitas
periódicas, toros invariantes y caos en sistemas cuadráticos sin equilibrio,
un oscilador Duffing memristivo, Rikitake, sistemas de cuatro y cinco
dimensiones y sincronización. Es un catálogo de modelos y antecedentes, no un
teorema general de ocultedad. Cada modelo que se importe debe volver a pasar
por el contrato algebraico, el contrato del integrador y el sondeo de cuenca
de este reporte.

La resolución de los archivos nombrados sólo por eISBN reveló cuatro textos
adicionales de alta prioridad:

\begin{itemize}
  \item \emph{Chaotic Systems with Multistability and Hidden Attractors}
        \cite{WangKuznetsovChen2021Multistability}: sistemas con equilibrios
        estables, sin equilibrios, con curvas o superficies de equilibrios
        (pp.\ 29--146); sistemas hipercaóticos, fraccionarios, memristivos y
        jerk con atractores ocultos (pp.\ 149--308); y multistabilidad
        simétrica o asimétrica (pp.\ 311--375).
  \item \emph{Systems with Hidden Attractors}
        \cite{PhamVolosKapitaniak2017Systems}: clasificación compacta por
        equilibrios estables (pp.\ 21--35), infinitos equilibrios (37--50),
        ausencia de equilibrios (51--63), sincronización (65--77) y
        realización circuital (79--102).
  \item \emph{Nonlinear Dynamics} \cite{DatserisParlitz2022Nonlinear}:
        dinámica continua, definición y medición de caos, bifurcaciones,
        entropía, dimensión fractal, coordenadas de retardo y análisis de
        series de tiempo (pp.\ 1--119). Su código Julia es un complemento
        natural de la reproducción con \juliads, pero una llamada de software
        no sustituye las hipótesis del teorema que implementa.
  \item \emph{Elements of Applied Bifurcation Theory}
        \cite{Kuznetsov2023Bifurcation}: equivalencia topológica y estabilidad
        estructural (pp.\ 47--87), bifurcaciones de equilibrios y ciclos
        (89--228), órbitas homoclínicas y heteroclínicas (229--279) y análisis
        numérico de bifurcaciones (569--655).
\end{itemize}

Estos hallazgos también afinan el portafolio de candidatos. Las familias sin
equilibrio, con equilibrios estables, con continuos de equilibrios,
fraccionarias, jerk y memristivas son \emph{fuentes de modelos}, no candidatos
promovidos automáticamente. Para cada capítulo seleccionado deben congelarse
ecuaciones, parámetros, orden, condición inicial publicada y criterio de
cuenca; después se aplica la criba de realización escalar/MIMO, regularidad
PWA, compatibilidad de memoria y disponibilidad de un oráculo numérico. Esta
secuencia conserva la diferencia entre «aparece en un libro de atractores
ocultos» y «fue reproducido bajo el contrato HAFO».

\subsection{Otros archivos relevantes y por qué no cumplen la misma función}

El inventario amplio mostró que la biblioteca es mayor de lo que reflejaba la
primera edición. Los siguientes archivos sí aportan piezas concretas, pero se
leen selectivamente:

\begin{longtable}{@{}L{5.25cm}L{4.1cm}L{4.65cm}@{}}
\caption{Archivos complementarios identificados en Drive.}
\label{tab:drive-complementarios}\\
\toprule
Archivo & Contenido útil & Uso recomendado \\
\midrule
\endfirsthead
\toprule
Archivo & Contenido útil & Uso recomendado \\
\midrule
\endhead
\path{978-3-030-91222-2.pdf} &
Cálculo, varias variables, series y EDO, pp.\ 243--760. &
Prerrequisito de consulta; no reemplaza un texto de dinámica. \\
\path{978-3-031-46320-4.pdf} &
Álgebra lineal y espectral, análisis complejo, Fourier y Laplace,
pp.\ 1--717. &
Resolventes, \(s^q\), transferencia y balance armónico. \\
\path{978-3-030-59472-5.pdf} &
Variable compleja, Fourier, Laplace y distribuciones, pp.\ 57--127. &
Banco de ejercicios durante función descriptiva y transferencia. \\
\path{978-3-031-01249-5.pdf} &
Sistemas dinámicos, formulación hamiltoniana y sistemas disipativos,
pp.\ 9--152. &
Segundo puente hacia caos; la parte cuántica queda como lectura futura. \\
\path{978-3-031-55460-5.pdf} &
Sistemas LTI, Laplace, realimentación y osciladores. &
Después de la metodología, para realización física y lectura de lazo. \\
\path{978-3-031-36364-1.pdf} &
Impedancia compleja, Laplace y electrónica no lineal. &
Consulta para circuitos de Chua; no es fuente de ocultedad. \\
\path{Introducing Fractals - Nigel Lesmoir-Gordon.pdf} &
Introducción visual a autosimilitud y dimensión fractal. &
Motivación posterior a la definición matemática; no demuestra caos ni Wada. \\
\path{[Marion,_Thornton]_Classical_Dynamics_of_Particles_and_Systems.pdf} &
Mecánica clásica lagrangiana y hamiltoniana. &
Referencia futura para sistemas conservativos, no prerrequisito de HAFO. \\
\path{978-1-4939-1360-2.pdf},
\path{978-1-4939-1366-4.pdf},
\path{978-3-030-04360-5.pdf},
\path{978-3-030-38585-9.pdf} y
\path{978-3-030-48028-8.pdf} &
Historia y cursos generales de física o mecánica. &
Biblioteca general: no se insertan en la ruta principal de la tesis. \\
\path{01_mechanics_lim.pdf}, \path{02_electromagnetism_lim.pdf},
\path{03_quantum_lim.pdf}, \path{05_modernphysics_lim.pdf} y
\path{5_ThermalPhysics_Reif_BPC.pdf} &
Física general por áreas. &
Consulta disciplinaria; no anteceden a Caputo, Lur'e o cuencas. \\
\bottomrule
\end{longtable}

En la subcarpeta \texttt{Cálculo vectorial y álgebra lineal} también se
confirmaron, entre otros,
\path{[Grossman] Algebra lineal 7 ed.pdf},
\texttt{Algebra-Lineal-y-sus-}\allowbreak
\texttt{Aplicaciones-3ra-Edición-}\allowbreak
\texttt{David-C.-Lay.pdf},
\texttt{Cálculo Vectorial (Calidad Mejorada) }\allowbreak
\texttt{(Jerrold E. Marsden, Anthony J. Tromba) }\allowbreak
\texttt{(Z-Library).pdf}
y varios solucionarios. Se elige un solo texto principal de álgebra y uno de
cálculo vectorial; los duplicados y solucionarios sirven para práctica, no
para crear otra etapa.

Esta clasificación es deliberadamente negativa cuando corresponde. Una guía
completa no consiste en recomendar todos los archivos: también debe explicar
qué libros pueden posponerse sin dejar un hueco matemático en la tesis.

\paragraph{Cierre del inventario que antes estaba pendiente.}
Los veintiún PDF de la lista provisional se abrieron uno por uno en Drive. El
visor permitió verificar portada, número total de páginas y esquema interno.
En \path{978-3-031-55202-1.pdf} el esquema no estaba embebido y en
\path{978-981-19-3273-1.pdf} el texto no era extraíble; en esos dos casos se
contrastó la identificación y la tabla de contenido con la ficha editorial
primaria de Springer. Se añadió además \path{978-3-030-65129-9.pdf}, visible en
la misma raíz. Por tanto, ninguno de estos nombres queda ya como ``pendiente''.

\begin{longtable}{@{}L{3.10cm}L{4.35cm}L{4.35cm}L{2.70cm}@{}}
\caption{Identificación cerrada de los archivos nombrados mediante eISBN.}
\label{tab:drive-isbn-cerrado}\\
\toprule
Archivo & Identificación verificada & Núcleo del índice & Decisión para HAFO \\
\midrule
\endfirsthead
\toprule
Archivo & Identificación verificada & Núcleo del índice & Decisión para HAFO \\
\midrule
\endhead
\path{978-3-030-63643-2.pdf} &
Volker Ziemann, \emph{Physics and Finance}. &
Procesos estocásticos, Fokker--Planck, series de tiempo, fractales y control
óptimo. & Consulta digital. \\
\path{978-3-030-87742-2.pdf} &
Samir Khene, \emph{Topics and Solved Exercises at the Boundary of Classical
and Modern Physics}. &
Átomos, interacción materia--radiación, cuerpo negro y termodinámica. &
Tangencial. \\
\path{978-3-031-35074-0.pdf} &
Francesco Scotognella, \emph{Exercises in Classical Physics---Mechanics and
Thermodynamics}. &
Cinemática, dinámica, energía, cuerpo rígido, fluidos y termodinámica. &
Repaso opcional. \\
\path{978-3-031-55202-1.pdf} &
José Rachid Mohallem, \emph{Lagrangian and Hamiltonian Mechanics}. &
Principio variacional, mecánica lagrangiana, hamiltoniana y teoría de campos. &
Consulta selectiva. \\
\path{978-3-031-67940-7.pdf} &
Teruo Matsushita, \emph{Exercises in Electricity and Magnetism}. &
Electrostática, corrientes, magnetismo, inducción, Maxwell y ondas. &
Apoyo circuital. \\
\path{978-3-031-82834-8.pdf} &
Teruo Matsushita, \emph{Electricity and Magnetism}, tercera edición. &
Fenómenos eléctricos y magnéticos estáticos y dependientes del tiempo. &
Apoyo circuital. \\
\path{978-3-031-82904-8.pdf} &
Tarek I. Zohdi, \emph{A Beginner's Primer: Basic Concepts in Mechanics}. &
Esfuerzo, deformación, vigas, torsión, elementos finitos y continuo. &
Fuera de ruta. \\
\path{978-3-031-84977-0.pdf} &
Vakhtang Putkaradze, \emph{A Concise Introduction to Classical Mechanics}. &
Variedades de configuración, estabilidad lineal, Hamilton, formas
diferenciales y transformaciones canónicas. & Alta selectiva. \\
\path{978-3-031-93666-1.pdf} &
Rudolf P. Huebener, \emph{Conductors, Semiconductors, Superconductors}, cuarta
edición. &
Bandas, semiconductores, superconductividad, magnetismo y nanoestructuras. &
Fuera de ruta. \\
\path{978-3-319-00750-2.pdf} &
Antonis Modinos, \emph{From Aristotle to Schrödinger}. &
Historia de la mecánica, termodinámica, electromagnetismo y cuántica. &
Lectura cultural. \\
\path{978-3-319-21816-8.pdf} &
Kerry Kuehn, \emph{A Student's Guide Through the Great Physics Texts}, vol.
III. &
Electricidad, magnetismo, óptica, energía y textos históricos. &
Lectura histórica. \\
\path{978-3-319-21828-1.pdf} &
Kerry Kuehn, \emph{A Student's Guide Through the Great Physics Texts}, vol.
IV. &
Calor, entropía, difusión, teoría atómica y cuántica. & Contexto remoto. \\
\path{978-3-319-51893-0.pdf} &
Juan C. Vallejo y Miguel A. F. Sanjuán, \emph{Predictability of Chaotic
Dynamics}. &
Exponentes de Lyapunov de tiempo finito, regímenes, escalas y predictibilidad. &
Alta prioridad. \\
\path{978-3-319-78361-1.pdf} &
Giovanni Landi y Alessandro Zampini, \emph{Linear Algebra and Analytic Geometry
for Physical Sciences}. &
Matrices, determinantes, sistemas, diagonalización y teoremas espectrales. &
Prerrequisito duplicado. \\
\path{978-3-0348-0697-8.pdf} &
B. Duplantier, S. Nonnenmacher y V. Rivasseau (eds.), \emph{Chaos: Poincaré
Seminar 2010}. &
Atractor de Lorenz, su topología y lectura estadística; dínamos y caos
cuántico. & Alta selectiva. \\
\path{978-3-662-46221-8.pdf} &
Bronshtein, Semendyayev, Musiol y Mühlig, \emph{Handbook of Mathematics}, sexta
edición. &
EDO, sistemas dinámicos y caos, transformadas y análisis numérico. &
Consulta digital. \\
\path{978-3-662-53685-8.pdf} &
Dietrich Stauffer, H. Eugene Stanley y Annick Lesne, \emph{From Newton to
Mandelbrot}, tercera edición. &
Fractales, sistemas dinámicos, bifurcaciones y caos, capítulos 5--6. &
Media--alta. \\
\path{978-981-15-2471-4.pdf} &
Rajendra K. Bera, \emph{The Amazing World of Quantum Computing}. &
Mecánica cuántica, Fourier, compuertas y algoritmos cuánticos. &
Fuera de ruta. \\
\path{978-981-16-7802-8.pdf} &
Kaushik Bhattacharya y Soumik Mukhopadhyay, \emph{Introduction to Advanced
Electrodynamics}. &
Fourier, covariancia, invariancia gauge y principio de acción. & Tangencial. \\
\path{978-981-16-8139-4.pdf} &
Amal Kumar Raychaudhuri, \emph{Classical Theory of Electricity and Magnetism}. &
Maxwell, ondas, plasma, magnetohidrodinámica y formulación variacional. &
Apoyo físico. \\
\path{978-981-19-3273-1.pdf} &
Kehui Sun, Shaobo He y Huihai Wang, \emph{Solution and Characteristic Analysis
of Fractional-Order Chaotic Systems}. &
Aproximación frecuencial, predictor--corrector, Adomian, algoritmos, dinámica,
complejidad y circuitos. & Prioridad máxima. \\
\path{978-3-030-65129-9.pdf} &
R. Prasad, \emph{Analog and Digital Electronic Circuits}. &
Análisis de circuitos, Laplace, Fourier, realimentación, diodos y
amplificadores. & Consulta circuital. \\
\bottomrule
\end{longtable}

La prioridad anterior se refiere a esta tesis, no a la calidad general del
libro. El núcleo añadido a la ruta queda formado por
\path{978-981-19-3273-1.pdf} y \path{978-3-319-51893-0.pdf}; la formación
geométrica se refuerza con \path{978-3-031-84977-0.pdf},
\path{978-3-0348-0697-8.pdf} y los capítulos 5--6 de
\path{978-3-662-53685-8.pdf}. Los restantes permanecen disponibles como
consulta digital o se excluyen explícitamente de la secuencia principal.

\subsection{Qué leer alrededor de cada procedimiento de tesis}
\label{sec:lecturas-por-procedimiento}

\begin{longtable}{@{}L{3.1cm}L{4.0cm}L{4.1cm}L{3.0cm}@{}}
\caption{Lectura antes--durante--después enlazada con el desarrollo del reporte.}
\label{tab:lectura-procedimientos}\\
\toprule
Procedimiento & Antes & Durante & Después \\
\midrule
\endfirsthead
\toprule
Procedimiento & Antes & Durante & Después \\
\midrule
\endhead
Equilibrios, jacobiano y Matignon &
Kantorovich I: varias variables y EDO; Hirsch--Smale--Devaney, caps.\ 7--9. &
\cref{sec:modelos,sec:fundamentos}: reducción escalar, pertenencia regional,
polinomio y margen angular. &
Comparar estabilidad local, índice y cuenca; resolver un caso PWA y uno suave. \\
Transferencia \(W_q\) &
Kantorovich II: álgebra espectral, compleja y Laplace; Das: ramas complejas. &
\cref{sec:lure_transferencia}: resolvente, menores, Cramer, rama de \(s^q\) y
convención de signo. &
Recuperar \(q=1\), verificar partes real e imaginaria y localizar polos. \\
Función descriptiva centrada y sesgada &
Cicogna: Fourier/Laplace; proyección ortogonal y saturación por intervalos. &
\cref{sec:lure_transferencia}: las integrales se derivan en el punto de uso. &
Comprobar límites \(A\to0\), \(A\to\infty\) y el residuo de cierre físico. \\
Machado &
Función descriptiva clásica, potencias complejas y selección de rama. &
\cref{sec:machado-fdf}: separar \(N^\mu\) auxiliar de la proyección física. &
Comparar semillas y ambos residuos; no promover una trayectoria por el cierre
auxiliar solamente. \\
Continuación causal Caputo &
Volterra, memoria completa; Kuznetsov para la idea de continuación y Amster
para la lectura topológica. &
\cref{sec:integracion_continuacion}: homotopía afín, historia única y reinicio
como operación distinta. &
Refinar paso, política de memoria y parámetro; estudiar operadores sobre
historias. \\
Puntos perpetuos &
Regla de cadena entera, jacobiano, flujo y cambios de coordenadas. &
\cref{sec:perpetual-points}: PP entero, obstrucciones Caputo, FPP--A y FPP--H. &
Pilotos con refinamiento de memoria; investigar formulación covariante o
funcional sin llamarla teoría establecida. \\
Chua entero y fraccionario &
Hirsch--Smale--Devaney, caps.\ 12 y 16; textos fraccionarios de la etapa 4. &
Capítulos de casos: sustituir parámetros en las identidades ya demostradas,
sin volver a esconder álgebra en apéndices. &
Pham et al., Pham--Volos--Kapitaniak y Wang--Kuznetsov--Chen para comparar
familias y realizaciones. \\
Kalman--Fitts, MAVPD y PLL &
Poincaré, Floquet, cilindro y cuencas; Hoppensteadt, \S\S2.2--4. &
Expedientes enteros: realización, continuación, regularidad y ocultedad. &
Kuznetsov: bifurcaciones; capítulos electrónicos sólo si se construye una
realización física. \\
Lorenz, Rabinovich--Fabrikant y candidatos &
Hirsch--Smale--Devaney, cap.\ 14; Caputo, Matignon y realización MIMO. &
Capítulos fraccionarios no Chua y \cref{sec:candidate-selection}. &
Ejecutar el protocolo completo antes de convertir un ejemplo bibliográfico en
resultado HAFO. \\
\bottomrule
\end{longtable}

\subsection{Correspondencia temática con la ruta única}

Los bloques siguientes no son una segunda secuencia. Traducen las etapas
numeradas de la \cref{tab:ruta-unica-biblioteca} a preguntas temáticas y
evitan mezclar prematuramente cálculo de operadores, geometría clásica y
dinámica con memoria.

\begin{itemize}[leftmargin=0pt,label={}]
 \item \textbf{Bloque A: base analítica.} Oliveira, funciones especiales,
 Mittag--Leffler y transformadas; reproducir cada identidad con cálculo
 simbólico y revisar dominios de convergencia.
 \item \textbf{Bloque B: derivadas e inicialización.} Comparar Riemann--Liouville y
 Caputo en Oliveira y Das; anotar terminal inferior, regularidad e información
 inicial requerida.
 \item \textbf{Bloque C: EDO geométrica.} Estudiar flujo, órbita, invariancia,
 linealización, mapa de Poincaré y conjuntos límite en
 Hirsch--Smale--Devaney y en los tomos geométricos de la UNAM.
 \item \textbf{Bloque D: existencia topológica.} Estudiar disparo, punto fijo, grado de
 Brouwer y Leray--Schauder con Amster; resolver después los ejemplos de la
 \cref{sec:metodos-topologicos-existencia}.
 \item \textbf{Bloque E: caos geométrico.} Relacionar Lorenz, dinámica simbólica,
 homoclinicidad y herradura con las secciones numéricas, sin identificar
 retrato complicado con prueba de caos.
 \item \textbf{Bloque F: memoria.} Leer no localidad en Herrmann e inicialización en
 Das, derivar la ecuación de Volterra y construir el estado funcional de la
 \cref{sec:geometria-topologia-caputo}.
 \item \textbf{Bloque G: retorno fraccionario.} Sustituir el mapa puntual de Poincaré
 por un operador de retorno sobre historias; estudiar cuándo su proyección al
 espacio \(x\) deja de ser un mapa univaluado.
 \item \textbf{Bloque H: aplicación HAFO.} Para cada caso, completar una ficha con
 espacio de fase, regularidad, simetría, equilibrios, conjuntos invariantes,
 cuenca, método de semilla y nivel de evidencia.
\end{itemize}

\subsection{Ejercicios de lectura activa y autocorrección}

\begin{enumerate}[label=\textbf{BL\arabic*.}]
 \item Demuestre \eqref{eq:library-order-semigroup} para funciones de prueba
 adecuadas usando integrales beta. Enumere las hipótesis que permiten cambiar
 el orden de integración.
 \item Para \(D_C^q x=\lambda x\), escriba la solución con
 Mittag--Leffler y compruebe numéricamente que, para \(0<q<1\), en general
 \(E_q(\lambda(t+s)^q)\ne E_q(\lambda t^q)E_q(\lambda s^q)\).
 \item Dibuje el dominio cortado usado para definir la rama principal de
 \(s^q\). Calcule los valores límite a ambos lados del corte negativo.
 \item En el PLL, identifique los bordes que deben pegarse para obtener
 \(\R\times S^1\). Explique qué error introduce medir la fase sin envolver.
 \item Construya el mapa de Poincaré de un ciclo plano sencillo. Liste las
 hipótesis que garantizan tiempo de retorno local continuo.
 \item Tome una figura caótica archivada de un expediente. Escriba tres afirmaciones que la
 figura apoya y tres que no demuestra: invariancia, transitividad, entropía
 positiva, homoclinicidad, atracción y ocultedad.
 \item Para una historia Caputo almacenada, construya dos perfiles con el mismo
 valor actual pero pasado distinto. Explique por qué la proyección del retorno
 puede ser multivaluada.
 \item Seleccione un capítulo del volumen de atractores ocultos y redacte el
 contrato mínimo para convertirlo en candidato HAFO: ecuaciones exactas,
 parámetros, orden, semilla, integrador, memoria, clasificación dinámica y
 prueba de cuenca.
\end{enumerate}

\paragraph{Pistas y criterios de comprobación.}
En BL1 debe aparecer la identidad beta y una justificación de Fubini. BL2 se
puede refutar ya con términos de orden \(ts\) o mediante valores de alta
precisión. BL3 debe mostrar la discontinuidad de argumento en \(\pm\pi\). BL4
debe usar una distancia cilíndrica. BL5 exige transversalidad al campo. BL6
debe separar evidencia finita de definición asintótica. BL7 queda incompleto
si sólo compara \(x(t)\); debe comparar el perfil de memoria. BL8 no puede
concluir ``oculto'' sólo porque el capítulo use esa palabra: debe reconstruir
el sondeo respecto de todos los equilibrios relevantes.
